\begin{lemma}[Strict shortening after the two-edge closure]
If \(C\) is a \(\noderef{TchangeCounterexample}\) and \(D\) is a
\(\noderef{TchangeKempeClosureData}\) for \(C\), then there is another
\(T\)-change counterexample \(C'\) whose
\(\noderef{TchangeCounterexampleMeasure}\) is strictly smaller than that
of \(C\).
\end{lemma}
\begin{proof}
Write \(U\) for the closed elementary tree and \(\sigma\) for the coloring
in \(D\).  For each \(i\), traverse the corresponding \((a,b)\)-walk from
the remote end \(a_i\), and stop at its first vertex \(c_i\) in \(V(U)\).
Such a vertex exists because \(b_i\in V(U)\), while \(a_i\notin V(U)\).
Call the resulting first-entry subpath \(L_i\), and let \(g_i\) be its
final exact edge.  Then \(L_i\) is nontrivial, simple, has only \(c_i\)
in \(U\), and its remote end still misses one of the two colors.  Thus it
satisfies every clause of \(\noderef{TreeExitPath}\) relative to \(U\).
The three \(L_i\) remain pairwise support-disjoint and retain their
distinct remote ends.

The two vertices adjacent to the edge-node \(g_i\) are the last outside
vertex and \(c_i\).  By \(\noderef{Incident}\) and the loopless endpoint
law in \(\noderef{LooplessMultigraph}\), these are the two distinct ends
of \(g_i\); hence
\(g_i\in\noderef{BoundaryEdges}(G,V(U))\).  Its color is \(a\) or \(b\).
But \(b\) is missing at \(b_3\in U\), and \(U\) is closed, so no boundary
edge has color \(b\), by \(\noderef{ColoredVertexSetClosed}\).
Therefore every \(g_i\) has color \(a\).
The vertices \(c_i\) are distinct: equality \(c_i=c_j\), together with
support-disjointness, would give two distinct \(a\)-colored exact edges
incident with one vertex, contrary to \(\noderef{AdjacentEdges}\) and
\(\noderef{IsProperEdgeColoring}\).

We verify that each \(L_i\) is the required prefix of a full component,
not merely an arbitrary exit walk.  At \(a_i\), one of \(a,b\) is missing.
Here missingness and presence are the complementary sets
\(\noderef{MissingColors}\) and \(\noderef{PresentColors}\).
There is at most one incident edge of the other color, since two such
exact identities would be adjacent and equally colored.  The nontrivial
path supplies one, so the incidence-graph degree of \(a_i\) is exactly
one.  The component contains an edge-node and a degree-one node, hence it
is a \(\noderef{TwoColorNontrivialPathComponent}\).  Apply
\(\noderef{TwoColorComponentSpanningPath}\) and orient the resulting
spanning path from \(a_i\).  The degree bound
\(\noderef{TwoColorIncidenceDegreeAtMostTwo}\) forces this full path and
\(L_i\) to agree successively from their common degree-one end through
\(c_i\).  Thus \(L_i\) is precisely the initial segment of its full
component up to the first vertex of \(U\).

The color \(b\) is missing on \(U\).  Since \(g_1\) is an \(a\)-colored
boundary identity, closedness shows that \(a\) cannot be missing at
any vertex of \(U\): otherwise that same boundary identity would violate
\(\noderef{ColoredVertexSetClosed}\).  Thus, after reversing the order
of the two colors, the new data have the convention that the first color
is missing on the tree and the second is not missing anywhere on it.
The \(k\)-triple, closed Tashkinov-tree,
and elementariness data are supplied by \(D\), with properness read from
\(\noderef{KTriple}\).  The three paths, their
disjointness and distinct ends, and the exact boundary witnesses \(g_i\)
in \(\noderef{BoundaryColorEdges}\) therefore form a new
\(\noderef{TchangeCounterexample}\) \(C'\) for the
color pair \((b,a)\).

Finally, both endpoints of \(f_1\) lie in \(U\).  The first-entry path
\(L_1\) stops before traversing \(f_1\), whereas the old first path contains
the edge-node \(f_1\).  Each \(L_i\) is a prefix of the corresponding old
path, so no length increases and the first length decreases strictly.
Consequently
\[
 |L_1|+|L_2|+|L_3|<|P_1|+|P_2|+|P_3|,
\]
which is exactly the asserted strict inequality of counterexample
measures.
\end{proof}
